% answers_short.tex —— 2004 年·北京卷·文科 · 简短答案速查
% 由 tools/make_short_answers.py 自 answers.tex 的 \gksolhead/\gkanswer 提取，
% 并与 2004/2004_数学真题_北京文_答案要点.md 逐题比对后生成，勿手改（改源 MD 后重跑）。

\begin{gknotice}
\textbf{使用说明}：本节只给最终答案便于快速核对，完整推导见解析版。答案为编写者独立推导，\textbf{非官方标准答案}。
\end{gknotice}

\gkpart{选择题}
\begin{center}\begin{tabular}{ccccccccc}
\toprule
题号 & 1 & 2 & 3 & 4 & 5 & 6 & 7 & 8 \\
\midrule
答案 & D & C & A & A & B & D & C & D \\
\bottomrule
\end{tabular}\end{center}

\gkpart{填空题}
\begin{itemize}[leftmargin=2.6em]
  \item[\textbf{9.}] $\pi$
  \item[\textbf{10.}] $x=1$ 或 $x=2$
  \item[\textbf{11.}] $4\sqrt{3}$，以及 $192\pi$
  \item[\textbf{12.}] $(0,-1)$，以及 $[1-\sqrt{2}, 1+\sqrt{2}]$
  \item[\textbf{13.}] 大，以及 $-3$
  \item[\textbf{14.}] $3$，以及 $52$（或计算式 $10\times 5+2$）
\end{itemize}

\gkpart{解答题}
\begin{itemize}[leftmargin=2.6em]
  \item[\textbf{15.}] $\tan A=-(2+\sqrt{3})$；$\triangle ABC$ 的面积为 $\dfrac{3(\sqrt{6}+\sqrt{2})}{4}$
  \item[\textbf{16.}] （1）$2\sqrt{10}$；（2）该最短路线的长为 $2\sqrt{5}$，$\dfrac{A_1M}{AM}=1$；（3）$\dfrac{\pi}{4}$（或 $45^\circ$）
  \item[\textbf{17.}] （1）该抛物线的方程为 $y^2=4x$，准线方程为 $x=-1$；（2）$y_1+y_2=-4$，直线 $AB$ 的斜率为 $-1$
  \item[\textbf{18.}] （1）$f(0)=0$，$f\left(\dfrac{1}{2}\right)=\dfrac{1}{2}$，$f\left(\dfrac{1}{4}\right)=\dfrac{1}{4}$，$f\left(\dfrac{1}{2^i}\right)=\dfrac{1}{2^i}$（$i=1, 2, \cdots$）；（2）$a_1=\dfrac{1}{4}$，$a_2=\dfrac{1}{16}$，$\lim_{n\to\infty}(a_1+a_2+\cdots+a_n)=\dfrac{1}{3}$
  \item[\textbf{19.}] （1）$B$ 站运行误差为 $\left|\dfrac{900}{v}-7\right|$ 分钟，$C$ 站运行误差为 $\left|\dfrac{1080}{v}-11\right|$ 分钟；（2）原题参数下两站运行误差之和恒大于等于 $\dfrac{13}{6}$ 分钟，满足条件的速度 $v$ 不存在（若原卷 $AB=15\text{ km}$ 属笔误实为 $5\text{ km}$，则解集为 $[39, 48.75]$）
  \item[\textbf{20.}] （1）$r_1\le r_2\le \cdots \le r_N$，除第 $N$ 组外每组至少含有 $3$ 个数；（2）余下的每个数都严格大于 $r_n$，证明见解析；（3）证明见解析
\end{itemize}
